\section{F7 · Den diskrete Fouriertransformen (DFT) og FFT}

\textbf{Kjernen:} DFT-en er Fourierkoeffisientene regnet ut fra $N$ punktverdier med trapesregelen, altså en matrise-vektor-multiplikasjon. Frekvenser som skiller seg med $N$, kan ikke skilles på gitteret (aliasing). FFT regner ut det samme på $O(N\log N)$ operasjoner.

\subsection{Fra Fourierkoeffisienter til DFT}
Trapesregelen på $[0,2\pi)$ med $N$ punkter $x_j=\frac{2\pi j}N$ gir
\[
c_k(f)=\frac1{2\pi}\int_0^{2\pi}f(x)\ee^{-\ii kx}\dd x\approx\frac1N\sum_{j=0}^{N-1}f\Bigl(\frac{2\pi j}{N}\Bigr)\ee^{-2\pi\ii jk/N}.
\]
Vi ser derfor på et signal $\vb f=(f_0,\dots,f_{N-1})^\T\in\C^N$ med $f_j=f(2\pi j/N)$.

\begin{formel}[DFT og IDFT\fs]
\[
\hat f_k=\frac1N\sum_{j=0}^{N-1}f_j\,\ee^{-2\pi\ii jk/N},\quad k=0,\dots,N-1,\qquad
f_j=\sum_{k=0}^{N-1}\hat f_k\,\ee^{2\pi\ii jk/N},\quad j=0,\dots,N-1.
\]
Med $w_N=\ee^{-2\pi\ii/N}$ og Fouriermatrisen $\FN=\bigl(w_N^{jk}\bigr)_{j,k=0}^{N-1}$:
\[
\hat{\vb f}=\tfrac1N\FN\vb f,\qquad \vb f=\overline{\FN}\,\hat{\vb f},\qquad \FN^{-1}=\tfrac1N\overline{\FN},\qquad \FN=\FN^\T.
\]
\end{formel}
\begin{teorem}[den sentrale summen (øving 8, oppg.~6)]
$\displaystyle\sum_{k=0}^{N-1}w_N^{jk}=\begin{cases}N,& j\equiv0\pmod N,\\0,&\text{ellers}\end{cases}$ \quad (geometrisk rekke med $q=w_N^j$, der $q^N=1$).
Det gir $\frac1N\overline{\FN}\FN=I$, altså at IDFT inverterer DFT.
\end{teorem}
\begin{eksempel}[små Fouriermatriser (kontrollregnet)]
\[
\mathcal F_2=\begin{pmatrix}1&1\\1&-1\end{pmatrix},\qquad
\mathcal F_4=\begin{pmatrix}1&1&1&1\\1&-\ii&-1&\ii\\1&-1&1&-1\\1&\ii&-1&-\ii\end{pmatrix}\quad(w_4=-\ii).
\]
For $N=5$ (Matte 4N vår 2023) skriver du $w_5^{jk}$ og reduserer eksponenten modulo 5: $(\mathcal F_5)_{jk}=w_5^{(jk\bmod 5)}$.
\end{eksempel}

\subsection{Egenskaper}
\begin{teorem}
\begin{itemize}
  \item \textbf{$N$-periodisk:} $\hat f_{k+N}=\hat f_k$ for alle $k\in\Z$, siden $w_N^N=1$. Indeksene $k=0,\dots,N-1$ (eller $-\frac N2,\dots,\frac N2-1$) er derfor nok.
  \item \textbf{Reelt signal:} $\hat f_{-k}=\overline{\hat f_k}$, altså $\hat f_{N-k}=\overline{\hat f_k}$. Spekteret er speilet, så bare halvparten trengs.
  \item \textbf{Aliasing-formelen:} er $f$ kontinuerlig med $\sum_k\abs{c_k(f)}<\infty$, så er
  \[
  \hat f_k=\sum_{\ell\in\Z}c_{k+\ell N}(f).
  \]
  På gitteret kan ikke frekvensene $k$ og $k+N$ skilles fra hverandre, fordi $\ee^{\ii(k+N)x_j}=\ee^{\ii kx_j}$.
\end{itemize}
\end{teorem}

\begin{teorem}[båndbegrensning, rekonstruksjon og Nyquist]
$f$ er \emph{båndbegrenset} med $B$ hvis $c_k(f)=0$ for $\abs k>B$. Hvis $N\ge2B+1$, er det høyst ett ledd i aliasing-summen for hver $k$, og
\[
c_k(f)=\hat f_k\ \ (0\le k\le B),\qquad c_{-k}(f)=\hat f_{N-k}\ \ (1\le k\le B),\qquad c_k(f)=0\ \text{ellers}.
\]
Da kan $f$ rekonstrueres eksakt fra prøvene (trigonometrisk interpolasjon). $2m+1$ prøver for høyeste frekvens $m$ kalles Nyquist-raten. Er $N\le 2B$, finnes det flere båndbegrensede funksjoner med de samme prøvene, og rekonstruksjonen er ikke entydig.
\end{teorem}
\begin{eksempel}[prøveeksamen, oppg.~7, kort]
$N=6$ og $\hat{\vb f}=(2,1,0,0,0,1)$. Med $B=2\le\frac{N-1}2$ blir $c_0=2$, $c_1=\hat f_1=1$ og $c_{-1}=\hat f_5=1$, så $f(x)=2+2\cos x$. Med $B=4$ er $N=6<2\cdot4+1$, og $f$ er ikke entydig.
\end{eksempel}

\begin{oppskrift}[tolke og filtrere et spekter (øving 8, oppg.~3)]
\begin{itemize}
  \item Indeks $k>\frac N2$ svarer til negativ frekvens $k-N$. \emph{fftshift} flytter dem slik at rekkefølgen blir $-\frac N2,\dots,\frac N2-1$.
  \item Med prøveavstand $\Delta t$ svarer indeks $k$ til frekvensen $\frac{k}{N\Delta t}$ Hz. Plott $\abs{\hat f_k}$ for $k=0,\dots,\lfloor N/2\rfloor$.
  \item \textbf{Filtrering eller støyfjerning:} sett $\hat f_k=0$ for uønskede frekvenser (eller der $\abs{\hat f_k}$ er liten) og ta IDFT. For et reelt signal må både $k$ og $N-k$ fjernes, ellers blir resultatet komplekst.
\end{itemize}
\end{oppskrift}

\subsection{Fast Fourier Transform (FFT)}
\begin{itemize}
  \item DFT som matrise-vektor-produkt koster $O(N^2)$. Med lyd i 44,1\,kHz er $N=44\,100$ per sekund, og det blir for dyrt.
  \item \textbf{Radix-2 (Cooley--Tukey, $N=2^n$, $M=N/2$):} del summen i første og andre halvdel, og se på like og odde $k$ hver for seg (uten faktoren $\frac1N$):
  \[
  \hat f_{2\kappa}=\sum_{j=0}^{M-1}\bigl(f_j+f_{j+M}\bigr)w_M^{j\kappa},\qquad
  \hat f_{2\kappa+1}=\sum_{j=0}^{M-1}\bigl(f_j-f_{j+M}\bigr)w_N^{j}\,w_M^{j\kappa}.
  \]
  Det gir to DFT-er av lengde $N/2$, $N$ addisjoner og $N/2$ multiplikasjoner. Gjentatt $n=\log_2N$ ganger gir det totalt $O(N\log N)$ operasjoner. Formlene er kontrollregnet.
  \item \textbf{Historie:} Gauss (ca.\ 1801, banen til Ceres, $N=12=4\cdot3$), Runge (1905) og Cooley--Tukey (1965). Andre faktoriseringer $N=N_1N_2$ fungerer også.
\end{itemize}

\subsection{Diskret sinus- og cosinustransform (DST og DCT)}
Samme idé som like og odde utvidelser. Et reelt signal gir reelle koeffisienter. Det finnes fire varianter av hver (I--IV), avhengig av hvordan signalet utvides. DCT-II er mest brukt, blant annet i bildekomprimering (JPEG). Å kaste små koeffisienter gir sterk komprimering, som i Barbara-eksempelet fra forelesningen (172\,kB ned til 3\,kB).

\begin{oving}
\begin{itemize}[leftmargin=5.6em,labelwidth=5em,labelsep=0.6em,align=left]
\item[\lett] \href{https://github.com/gustavbv/Matte4C---eksamen/blob/main/oppgavesett/F07-dft-og-fft/2-Bergmann-eksamener/010-TMA4125-2022-var-oppg7-oppgave.pdf}{Matte 4N vår 2022, oppg.~7} {\color{graa}(\href{https://github.com/gustavbv/Matte4C---eksamen/blob/main/oppgavesett/F07-dft-og-fft/2-Bergmann-eksamener/010-TMA4125-2022-var-oppg7-svar.pdf}{LF})}: $\mathcal F_4$, DFT og invers DFT
\item[\lett] \href{https://github.com/gustavbv/Matte4C---eksamen/blob/main/oppgavesett/F07-dft-og-fft/2-Bergmann-eksamener/008-TMA4125-30-2023-kont-oppg3-oppgave.pdf}{Matte 4N kont.\ 2023, oppg.~3} {\color{graa}(\href{https://github.com/gustavbv/Matte4C---eksamen/blob/main/oppgavesett/F07-dft-og-fft/2-Bergmann-eksamener/008-TMA4125-30-2023-kont-oppg3-svar.pdf}{LF})}: DFT av samplet $\ee^{3\ii x}$
\item[\middels] \href{https://github.com/gustavbv/Matte4C---eksamen/blob/main/oppgavesett/F07-dft-og-fft/1-TMA4432-ovinger-og-proveeksamen/001-TMA4432-2026-proveeksamen-oppg7-oppgave.pdf}{Prøveeksamen, oppg.~7} {\color{graa}(\href{https://github.com/gustavbv/Matte4C---eksamen/blob/main/oppgavesett/F07-dft-og-fft/1-TMA4432-ovinger-og-proveeksamen/001-TMA4432-2026-proveeksamen-oppg7-svar.pdf}{LF})}: båndbegrensning og rekonstruksjon
\item[\middels] \href{https://github.com/gustavbv/Matte4C---eksamen/blob/main/oppgavesett/F07-dft-og-fft/2-Bergmann-eksamener/007-TMA4125-2023-var-oppg4-oppgave.pdf}{Matte 4N vår 2023, oppg.~4} {\color{graa}(\href{https://github.com/gustavbv/Matte4C---eksamen/blob/main/oppgavesett/F07-dft-og-fft/2-Bergmann-eksamener/007-TMA4125-2023-var-oppg4-svar.pdf}{LF})}: $\mathcal F_5$, båndbegrensning og aliasing
\item[\middels] \href{https://github.com/gustavbv/Matte4C---eksamen/blob/main/oppgavesett/F07-dft-og-fft/2-Bergmann-eksamener/005-TMA4125-30-2025-kont-oppg5-oppgave.pdf}{Matte 4N kont.\ 2025, oppg.~5} {\color{graa}(\href{https://github.com/gustavbv/Matte4C---eksamen/blob/main/oppgavesett/F07-dft-og-fft/2-Bergmann-eksamener/005-TMA4125-30-2025-kont-oppg5-svar.pdf}{LF})}: aliasing og Nyquist for $3+2\cos 2x$
\item[\vanskelig] \href{https://github.com/gustavbv/Matte4C---eksamen/blob/main/oppgavesett/F07-dft-og-fft/1b-TMA4432-ekstraoppgaver/003-TMA4432-oving-8-oppg6-oppgave.pdf}{Øving 8, oppg.~6 (ekstra)} {\color{graa}(\href{https://github.com/gustavbv/Matte4C---eksamen/blob/main/oppgavesett/F07-dft-og-fft/1b-TMA4432-ekstraoppgaver/003-TMA4432-oving-8-oppg6-svar.pdf}{LF})}: bevis: periodisitet og summen av $w_N^{jk}$
\item[\vanskelig] \href{https://github.com/gustavbv/Matte4C---eksamen/blob/main/oppgavesett/F07-dft-og-fft/3-hele-oppgaven-i-pensum/012-TMA4130-35-2023-host-oppg7-oppgave.pdf}{Matte 4N/4D høst 2023, oppg.~7} {\color{graa}(\href{https://github.com/gustavbv/Matte4C---eksamen/blob/main/oppgavesett/F07-dft-og-fft/3-hele-oppgaven-i-pensum/012-TMA4130-35-2023-host-oppg7-svar.pdf}{LF})}: tolke et spekter og kode for et filter
\end{itemize}
\end{oving}

